\begin{lemma}[The frozen operator satisfies (Q1)--(Q2)]
Let $n\ge2$, $p>1$, $0\le\varepsilon\le1$, $0<\lambda\le\Lambda$, $a>0$ and
$0<\nu_0<\nu_1$.  Let $\Omega\subseteq\R^n$ be open, let $f$ satisfy
(F1) of the thesis, (F2) of the thesis and (F3) of the thesis on $\Omega$ with parameters
$p,\varepsilon,\lambda,\Lambda,a,\nu_0,\nu_1$ and modulus $\sigma$
(\noderef{VariationalStructureF123}), let $x_1\in\Omega$ and let $z_1\in\R$.
Write $f^0(h)=f(x_1,z_1,h)$ and $A^0=\nabla_hf^0$
(\noderef{FrozenGradientOperator}).  Then the normalised frozen operator
\[
\widetilde A=A^0-A^0(0)
\]
satisfies (Q1) of the thesis--(Q2) of the thesis on all of $\R^n$ with the parameters
$p,\varepsilon$ and the structure constants $\gamma_0=\nu_0$,
$\gamma_1=\nu_1$ (\noderef{EllipticStructureQ12}), the field being read as the
field $(x,z,h)\mapsto\widetilde A(h)$, constant in the point and in the value.

\noindent\textbf{What this node is, and why it is a node.}  It is the step by
which \S4 and \S5 hand their variational data to the nodes of \S3.  Every
statement of \S3 that \S4 applies to a comparison map --- Lemma~Lemma~3.1 of the thesis,
Theorem~Theorem~1 of the thesis and Theorem~Theorem~2 of the thesis --- is stated for a field satisfying
(Q1) of the thesis--(Q2) of the thesis, and the field to which \S4 applies them is the
frozen operator of the density, normalised as the thesis prescribes at l.~542
(``Since $\operatorname{div}A(0)=0$ we may assume without loss of generality
that $A(0)=0$ and that $A$ satisfies conditions (Q1) of the thesis and
(Q2) of the thesis'').  That the frozen operator \emph{does} satisfy
(Q1) of the thesis--(Q2) of the thesis, with the constants $\nu_0,\nu_1$ of (F2) of the thesis
and not merely with some constants, is exactly what the thesis means by its
sentence at l.~1626 about the \S5 functional being ``of the type considered in
\S4''.  In \S5 the corresponding step needs no node, because there the data are
already given as (Q1) of the thesis--(Q2) of the thesis and only the a.e.\ quantifier has
to be upgraded.

\noindent\textbf{On the constants and the normalisation.}  No constant is
produced: $\gamma_0$ and $\gamma_1$ are the $\nu_0$ and $\nu_1$ of
(F2) of the thesis themselves, undegraded, so no quantifier-order question arises.
The normalisation $\widetilde A(0)=0$ is a definitional tautology and is
therefore not stated as a clause; what matters is that it costs nothing in the
structure conditions, since (Q1) of the thesis and (Q2) of the thesis constrain only the
derivative of the field and the derivative is unchanged by subtracting a
constant vector.  The restriction $\varepsilon\le1$ is the range in which \S4
and \S5 work, carried here for uniformity with the statements that consume this
node; the proof does not use the upper bound, and \S3 itself drops it (thesis
l.~543--545, ``in this section the restriction $\varepsilon\le1$ plays no
role'').
\end{lemma}
\begin{proof}
Write
\[
f^0(h)=f(x_1,z_1,h),\qquad
A^0=\nabla_hf^0\quad(\noderef{FrozenGradientOperator}),\qquad
\widetilde A=A^0-A^0(0),
\]
so that the field whose structure is asserted is $\widetilde A$, viewed as a
field on $\R^n$ that is constant in the point and in the value.  We verify the
six clauses of (Q1) of the thesis--(Q2) of the thesis for it on all of $\R^n$ with the
parameters $p,\varepsilon,\gamma_0=\nu_0,\gamma_1=\nu_1$
(\noderef{EllipticStructureQ12}).  Note first that subtracting the constant vector
$A^0(0)$ changes neither continuity, nor differentiability, nor any derivative
of the field, so it suffices to treat $A^0$; and that $A^0$ does not depend on
the point of $\R^n$, so the clauses that are required at every point of the
domain, or at almost every point, follow from the single computation below.
Write $w(h)=\varepsilon+|h|^2$ and $D=\{h:w(h)\ne0\}$, an open set.
\begin{itemize}
\item \emph{Carath\'eodory measurability in the point.}  For fixed $z$ and $h$
  the map $x\mapsto\widetilde A(x,z,h)$ is \emph{constant}, with value
  $A^0(h)-A^0(0)$, since $\widetilde A$ does not depend on its first argument;
  a constant function is Lebesgue measurable.
\item \emph{Carath\'eodory joint continuity in $(z,h)$.}  For fixed $x$ the map
  $(z,h)\mapsto\widetilde A(x,z,h)=A^0(h)-A^0(0)$ does not depend on $z$ either,
  so it is the composition of $h\mapsto A^0(h)-A^0(0)$ with the continuous
  projection $(z,h)\mapsto h$; and $A^0$ is continuous by the next clause.  A
  composition of continuous maps is continuous, so the joint map is continuous
  on $\R\times\R^n$.
\item \emph{Continuity.}  By (F2) of the thesis
  (\noderef{VariationalStructureF123}, clause \texttt{F2\_diff}, applied at the
  point $x_1\in\Omega$, which is a hypothesis, and at the value $z_1$) the function $f^0$ is of class $C^1$ on $\R^n$, so
  $h\mapsto\nabla_hf^0(h)=A^0(h)$ is continuous.
\item \emph{$C^1$ on $D$.}  By clause \texttt{F2\_diff2} of the same node,
  $f^0$ is of class $C^2$ on $D$; hence its gradient $A^0$ is of class $C^1$
  there.
\item \emph{(Q1) of the thesis.}  Fix $h\in D$ and $\xi\in\R^n$.  The $i$-th
  component of $A^0$ is $A^{0,i}(h)=D_{h_i}f^0(h)$, so the derivative of $A^0$
  at $h$ in the direction $\xi$ has $i$-th component $D_\xi D_{h_i}f^0(h)$, and
  therefore
  \[
  \sum_{i=1}^n\bigl(D_\xi A^0(h)\bigr)^i\xi^i
  =\sum_{i=1}^n\xi^iD_\xi D_{h_i}f^0(h)
  =D_\xi\Bigl(\sum_{i=1}^n\xi^iD_{h_i}f^0\Bigr)(h)
  =D_\xi\bigl(D_\xi f^0\bigr)(h),
  \tag{1}
  \]
  the middle equality because $\xi$ is a constant vector and differentiation in
  a fixed direction is linear.  The right-hand side is exactly the quantity
  that (F2) of the thesis estimates, and its lower bound there is
  $\nu_0w(h)^{\frac p2-1}|\xi|^2$, which is (Q1) of the thesis with
  $\gamma_0=\nu_0$.
\item \emph{(Q2) of the thesis.}  For $h\in D$ define the bilinear form
  $B(\zeta,\xi)=D_\xi(D_\zeta f^0)(h)$ on $\R^n\times\R^n$; it is bilinear
  because $f^0$ is twice differentiable at $h$ (clause \texttt{F2\_diff2} and
  openness of $D$).  By the symmetry of the second derivative (Clairaut,
  Preamble), whose hypothesis is exactly twice differentiability at the point,
  $B$ is symmetric.  By (F2) of the thesis and $(1)$,
  \[
  \nu_0w(h)^{\frac p2-1}|\zeta|^2\le B(\zeta,\zeta)\le
  \nu_1w(h)^{\frac p2-1}|\zeta|^2
  \qquad(\zeta\in\R^n),
  \tag{2}
  \]
  so $B$ is nonnegative.  For a symmetric nonnegative bilinear form,
  $|B(\zeta,\xi)|\le B(\zeta,\zeta)^{1/2}B(\xi,\xi)^{1/2}$: the quadratic
  polynomial $t\mapsto B(\zeta+t\xi,\zeta+t\xi)
  =B(\zeta,\zeta)+2tB(\zeta,\xi)+t^2B(\xi,\xi)$ is nonnegative for every real
  $t$, so its discriminant $4B(\zeta,\xi)^2-4B(\zeta,\zeta)B(\xi,\xi)$ is at
  most $0$ when $B(\xi,\xi)>0$, and when $B(\xi,\xi)=0$ the same nonnegativity
  for large $|t|$ of the affine function $B(\zeta,\zeta)+2tB(\zeta,\xi)$ forces
  $B(\zeta,\xi)=0$.  Taking $\zeta=e_i$ and $\xi=e_j$ the two standard basis
  vectors, which have $|e_i|=|e_j|=1$, and using the upper bound in $(2)$ twice,
  \[
  \bigl|\bigl(D_{e_j}A^0(h)\bigr)^i\bigr|=|B(e_i,e_j)|
  \le B(e_i,e_i)^{1/2}B(e_j,e_j)^{1/2}
  \le\bigl(\nu_1w(h)^{\frac p2-1}\bigr)^{1/2}\bigl(\nu_1w(h)^{\frac p2-1}\bigr)^{1/2}
  =\nu_1w(h)^{\frac p2-1},
  \]
  the first equality by the component computation of the previous item --- there
  $(D_{e_j}A^0(h))^i=D_{e_j}D_{h_i}f^0(h)=B(e_i,e_j)$ --- and the last step
  because $w(h)^{\frac p2-1}\ge0$ and $\nu_1>0$.  This is (Q2) of the thesis with
  $\gamma_1=\nu_1$.
\end{itemize}
Since the last four clauses hold for every point of $\R^n$ and in particular almost
everywhere, $\widetilde A$ satisfies (Q1) of the thesis--(Q2) of the thesis on $\R^n$ ---
and a fortiori on every subset of $\R^n$ --- with the parameters
$p,\varepsilon,\nu_0,\nu_1$, which is the conclusion.

\end{proof}
